\begin{pastexam}[label={px:df3-107-ds-05}]{地方三等資訊處理 107 年 第五題}
（申論；20 分）

斐波納契數（Fibonacci number）$F_n$ 的定義是 $F_0 = 0, F_1 = 1, F_n = F_{n-1} + F_{n-2}, n > 1$。計算 Fibonacci number $F_n$ 的演算法，以類似 C 語言表示如下：
\begin{quote}
\begin{tabular}{@{}r@{\qquad}l@{}}
1 & integer $f[N]$;\quad // array of $N$ integers \\
2 & integer $F(n)$ \{ \\
3 & \qquad \textbf{if} $(f[n] < 0)$ \\
4 & \qquad\qquad $f[n] = F(n-1) + F(n-2)$; \\
5 & \qquad return $f[n]$; \\
6 & \} \\
7 & integer $\mathrm{Fib}(n)$ \{ \\
8 & \qquad $f[0] = 0$; $f[1] = 1$; \\
9 & \qquad \textbf{for} $(i = 2;\ i \le n;\ i = i + 1)$ \\
10 & \qquad\qquad $f[i] = -1$; \\
11 & \qquad return $F(n)$; \\
12 & \}
\end{tabular}
\end{quote}
其中資料型態 integer 表示整數。假設輸入的整數 $n > 1$。主程式執行 $\mathrm{Fib}(n)$，則副程式 $F(n)$ 第 4 行之指令：

$f[n] = F(n-1) + F(n-2)$ 會被執行幾次？請說明理由。（20 分）
\end{pastexam}

\begin{pastsoln}
答案：$n - 1$ 次。 這是由上而下的記憶化（memoization，資結書第~18.1.2~節）：$f[k] = -1$ 表示「還沒算」，算過就存起來。
陣列大小 $N$ 要大於 $n$。

\textbf{理由}。 對每個 $k$（$0 \le k \le n$），看第 4 行在 $F(k)$ 裡執行幾次：
\begin{itemize}
\item $k = 0, 1$：第 8 行已設 $f[0] = 0$、$f[1] = 1$，都不小於 $0$，第 3 行的條件不成立，第 4 行\textbf{從不}在 $F(0)$、$F(1)$ 中執行。
\item $2 \le k \le n$：第 10 行把 $f[k]$ 設為 $-1$。 \textbf{第一次}呼叫 $F(k)$ 時 $f[k] = -1 < 0$，第 4 行執行一次，把 $f[k]$ 改成
$F_k \ge 1$。 之後再呼叫 $F(k)$，$f[k] \ge 1$，條件不成立，直接傳回。 所以只要 $F(k)$ 至少被呼叫一次，第 4 行在 $F(k)$ 裡\textbf{恰好}
執行一次。
\item 每個 $F(k)$（$2 \le k \le n$）都會被呼叫：$F(n)$ 由第 11 行呼叫；若 $F(k)$（$k \ge 3$）第一次被呼叫，它執行第 4 行時會先呼叫
$F(k-1)$。 由 $k = n$ 往下歸納，$F(n), F(n-1), \ldots, F(2)$ 都會被呼叫。
\end{itemize}
所以第 4 行一共執行 $n - 1$ 次（$k = 2, 3, \ldots, n$ 各一次）。

\textbf{追蹤 $n = 5$}（第 4 行先算 $F(n-1)$ 再算 $F(n-2)$，C 語言沒規定加法兩邊的計算順序，但不影響次數）：
\begin{center}
\small
\begin{tabular}{lll}
\toprule
呼叫 & $f$ 的值 & 動作 \\
\midrule
$F(5)$ & $f[5] = -1$ & 執行第 4 行（第 1 次），呼叫 $F(4)$ \\
\quad $F(4)$ & $f[4] = -1$ & 執行第 4 行（第 2 次），呼叫 $F(3)$ \\
\qquad $F(3)$ & $f[3] = -1$ & 執行第 4 行（第 3 次），呼叫 $F(2)$ \\
\qquad\quad $F(2)$ & $f[2] = -1$ & 執行第 4 行（第 4 次），$F(1) + F(0) = 1$，$f[2] = 1$ \\
\qquad\quad $F(1)$ & $f[1] = 1$ & 直接傳回 $1$；$f[3] = 1 + 1 = 2$ \\
\quad\quad $F(2)$ & $f[2] = 1$ & 直接傳回 $1$；$f[4] = 2 + 1 = 3$ \\
\quad $F(3)$ & $f[3] = 2$ & 直接傳回 $2$；$f[5] = 3 + 2 = 5$ \\
\bottomrule
\end{tabular}
\end{center}
第 4 行執行 $4 = 5 - 1$ 次，傳回 $F_5 = 5$。 順帶一提，$F$ 總共被呼叫 $1 + 2(n-1) = 2n - 1$ 次（每次執行第 4 行產生兩次呼叫），
所以 $\mathrm{Fib}(n)$ 的時間是 $\Theta(n)$；不記憶的直接遞迴要 $\Theta(\phi^{n})$（見 歷屆試題~\ref{px:gk3-096-ds-03} ㈢）。
\end{pastsoln}
